\chapter{Когда искать, а когда решать: добавляем \nt{НОРА}}
\chaptermark{Добавляем \nt{НОРА}}
\label{sec:nora}

\section{Чего не хватает}

В главе~\ref{sec:dofa} сеть училась благодаря шуму: случайные отклонения активности были поиском, а \nt{ДОФА} сообщал, стало ли от них лучше. Но там же остался незаданным один параметр --- \emph{сколько} искать. Если шум слишком силён, сеть мечется и не пользуется тем, что уже знает. Если слишком слаб, она раз за разом выбирает одно и то же и не замечает, что где-то рядом стало лучше. В обучении с подкреплением это называют выбором между использованием и поиском, и у каждого алгоритма для него есть ручка. В главе~\ref{sec:neurontypes} мы предположили, что в мозге эту ручку крутит \nt{НОРА}.

Норадреналиновых нейронов у человека несколько десятков тысяч (таблица~\ref{tab:types}), почти все они собраны в одной маленькой группе, а их выходы расходятся практически по всему мозгу, хотя разные части этой группы, как выяснилось, отчасти обслуживают разные области~\cite{Poe2020}. Это широковещательный канал, ещё более узкий, чем \nt{ДОФА}. Работает он в двух режимах~\cite{AstonJones2005}. В \emph{фоновом} нейроны выдают импульсы ровно, с частотой от долей до нескольких импульсов в секунду, и эта частота медленно меняется вместе с бодрствованием. В \emph{отрывистом} они отвечают короткой пачкой на важное для текущей задачи событие, примерно в момент принятия решения. Удобная, но неточная контрольная точка --- диаметр зрачка: он меняется вместе с активностью этих нейронов, но также вместе с активностью других областей~\cite{Joshi2016} и холинергических выходов в коре~\cite{Reimer2016}. Зрачок показывает общий уровень возбуждения, а не одну \nt{НОРА}.

\section{Коэффициент усиления}

Измерено вот что: норадреналин подавляет фоновую активность нейронов сильнее, чем их ответ на стимул, и отношение сигнала к фону растёт~\cite{Foote1975NE}. Что крутизна характеристики отдельного нейрона при этом буквально умножается на коэффициент, из опыта не следует. Как и в главе~\ref{sec:neurontypes}, возьмём простую модель, которая передаёт этот эффект: норадреналин меняет крутизну передаточной характеристики получателя~\cite{ServanSchreiber1990}. Слабые, подпороговые входы (фон) тогда дают ещё меньше, а сильные --- ещё больше. Пусть нейрон отвечает частотой
\begin{empheq}[box=\keybox]{equation}
r \;=\; F\bigl(g\,(u-\theta)\bigr), \qquad F(z)=\frac{r_{\max}}{1+e^{-z}} ,
\label{eq:gain}
\end{empheq}
где $u$ --- входной ток, $\theta$ --- порог, а $g$ --- коэффициент усиления, которым в модели управляет \nt{НОРА}. При малом $g$ частота плавно следует за входом в широком диапазоне: нейрон --- \emph{аналоговый усилитель}. При большом $g$ почти любой вход выше порога даёт полную частоту, а ниже --- ноль: нейрон --- \emph{компаратор}, почти цифровой элемент (рис.~\ref{fig:gaincurves}). Одна ручка переключает тот же нейрон между двумя режимами, которые в электронике делают разные схемы.

\begin{figure}[t]
\centering
\begin{tikzpicture}[>=Stealth,x=1.1cm,y=2.4cm]
\draw[->,gray!70!black] (-3.3,0) -- (3.4,0) node[below left,black] {\small $u$};
\draw[->,gray!70!black] (-3.3,0) -- (-3.3,1.2) node[left,black] {\small $r$};
\draw[densely dashed,gray!60!black] (0,0) -- (0,1.1);
\node[below,font=\small] at (0,0) {$\theta$};
\draw[densely dotted,gray!60!black] (-3.3,1) -- (3.2,1) node[right,black,font=\small] {$r_{\max}$};
\draw[very thick,blue!60!black] plot[domain=-3.2:3.2,samples=60] (\x,{1/(1+exp(-0.8*\x))});
\draw[very thick,orange!85!black] plot[domain=-3.2:3.2,samples=60] (\x,{1/(1+exp(-2*\x))});
\draw[very thick,red!70!black] plot[domain=-3.2:3.2,samples=120] (\x,{1/(1+exp(min(-8*\x,9)))});
\node[blue!60!black,font=\small,anchor=west] at (0.5,0.12) {малое $g$: усилитель};
\node[red!70!black,font=\small,anchor=east] at (-0.3,0.85) {большое $g$: компаратор};
\end{tikzpicture}
\caption{Передаточная характеристика~\eqref{eq:gain} при трёх значениях коэффициента усиления $g$.}
\label{fig:gaincurves}
\end{figure}

Помогает ли большое усиление против шума? Не всегда. Пусть два входа отличаются на $\Delta u$, и надо сказать, какой больше. Если шум $\sigma_{\text{до}}$ добавлен к входу \emph{до} усилителя, то усиливается и сигнал, и шум, и их отношение не меняется. Если же шум $\sigma_{\text{после}}$ добавляется \emph{после} усилителя --- от ненадёжных синапсов и случайных импульсов самой сети, как в главе~\ref{sec:code}, --- то отношение сигнала к шуму
\begin{empheq}[box=\keybox]{equation}
d' \;=\; \frac{g\,\Delta u}{\sqrt{g^2\sigma_{\text{до}}^2+\sigma_{\text{после}}^2}}
\label{eq:gainsnr}
\end{empheq}
растёт с $g$, пока $g\sigma_{\text{до}}$ не сравняется с $\sigma_{\text{после}}$~\cite{ServanSchreiber1990}.

\begin{keyblock}
\begin{proposition}[Когда усиление помогает]
\label{prop:gainnoise}
Увеличивая коэффициент усиления, \nt{НОРА} улучшает различение входов только против того шума, который возникает \emph{после} усилителя, в самой сети. Шум, пришедший вместе со входом, усиление не уберёт.
\end{proposition}
\end{keyblock}

\section{Усиление в выборе из двух}

Вернёмся к выбору из главы~\ref{sec:gaba}: две группы \nt{ГЛУТ}-нейронов со входами $I_1$, $I_2$ тормозят друг друга с силой $\beta$. Теперь у каждой группы есть коэффициент усиления $g$: в уравнениях выбора~\eqref{eq:wta} выражение в скобках умножается на $g$. Пока работают обе группы, установившиеся частоты находятся из $r_1=g(I_1-\beta r_2)$, $r_2=g(I_2-\beta r_1)$. Сложив и вычтя эти равенства, получим сумму $r_1+r_2=g(I_1+I_2)/(1+g\beta)$ и разность $r_1-r_2=g(I_1-I_2)/(1-g\beta)$. Их отношение --- насколько резко сеть предпочитает один вход другому:
\begin{empheq}[box=\keybox]{equation}
\frac{r_1-r_2}{r_1+r_2} \;=\; \varepsilon\,\frac{1+g\beta}{1-g\beta}, \qquad \varepsilon=\frac{I_1-I_2}{I_1+I_2}, \quad g\beta<1 .
\label{eq:wtagain}
\end{empheq}
Малая разница входов $\varepsilon$ усиливается в $(1+g\beta)/(1-g\beta)$ раз, и этот множитель неограниченно растёт, когда $g\beta$ подходит к единице. При $g\beta>1$ состояние, где работают обе группы, неустойчиво, как в утверждении~\ref{prop:wta}: побеждает одна, другая молчит (рис.~\ref{fig:pitchfork}).

\begin{figure}[t]
\centering
\begin{tikzpicture}[>=Stealth,x=3.2cm,y=1.5cm]
\draw[->,gray!70!black] (0,0) -- (2.15,0) node[below left,black] {\small $g\beta$};
\draw[->,gray!70!black] (0,-1.25) -- (0,1.35) node[above,black] {\small $(r_1-r_2)/(r_1+r_2)$};
\draw[gray!60!black] (1,-0.04) -- (1,0.04); \node[below right,font=\small] at (1,0) {$1$};
\node[left,font=\scriptsize] at (0,1) {$1$}; \node[left,font=\scriptsize] at (0,-1) {$-1$};
\draw[very thick,blue!60!black] plot[domain=0:0.905,samples=60] (\x,{0.05*(1+\x)/(1-\x)}) -- (1,1) -- (2.05,1);
\draw[very thick,blue!60!black] (1.105,-1) -- (2.05,-1);
\draw[thick,densely dashed,gray!60!black] (1,0) -- (2.05,0);
\node[font=\small,align=center] at (0.45,-0.62) {размышляет:\\обе группы\\работают};
\node[font=\small,align=center] at (1.55,0.45) {решила:\\работает одна};
\node[font=\small,blue!60!black,anchor=south west] at (1.05,-0.97) {слабый вход победил раньше --- держится};
\end{tikzpicture}
\caption{Выбор из двух при разном коэффициенте усиления, входы отличаются на $\varepsilon=0.05$. При $g\beta>1$ устойчивы оба решения (сплошные линии; победа слабого входа --- при $g\beta>(1+\varepsilon)/(1-\varepsilon)\approx1.1$), и сеть держит то, которое уже приняла; симметричное состояние (штриховая линия) неустойчиво.}
\label{fig:pitchfork}
\end{figure}

\begin{keyblock}
\begin{proposition}[Усиление переключает режим выбора]
\label{prop:gainwta}
В сети выбора с взаимным торможением $\beta$ коэффициент усиления $g$ переводит сеть через границу $g\beta=1$. Ниже неё сеть «размышляет»: обе группы работают, а разница входов усиливается по~\eqref{eq:wtagain}. Выше --- сеть «решила»: работает одна группа, и решение держится. Меняя $g$, \nt{НОРА} может переключать сеть между этими режимами, не трогая ни одного веса.
\end{proposition}
\end{keyblock}

\section{Усиление как температура}

Теперь добавим к выбору шум, возникающий в самой сети, после усилителя. Какая группа победит, решает знак величины $g(u_A-u_B)+\eta$, где $u_A-u_B$ --- разница входов, то есть разница выученных весов из главы~\ref{sec:dofa}, а $\eta$ --- шум с масштабом $\sigma$. Если шум распределён по логистическому закону (для гауссова кривая почти та же), вероятность выбрать $A$ равна
\begin{empheq}[box=\keybox]{equation}
P(A) \;=\; \frac{1}{1+e^{-g\,(u_A-u_B)/\sigma}} .
\label{eq:softmax}
\end{empheq}
Программист узнает здесь выбор по softmax с температурой $T=\sigma/g$. Коэффициент усиления --- обратная температура. При малом $g$ даже заметно лучший вариант выбирается ненамного чаще худшего: сеть много ищет. При большом $g$ лучший известный вариант выбирается почти всегда: сеть пользуется тем, что знает.

Уточним, что такое $g$ в~\eqref{eq:wtagain} и~\eqref{eq:softmax}: это \emph{действующее} усиление разницы входов текущей задачи в момент выбора. Два режима \nt{НОРА} действуют на него противоположно. Отрывистая пачка в момент решения повышает его, и сеть закрепляет выбор. Высокая фоновая активность, наоборот, идёт вместе с поиском: у людей перед выбором наугад базовый зрачок шире~\cite{JepmaNieuwenhuis2011}, а у крыс вход \nt{НОРА} в переднюю поясную кору переводит выбор в случайный режим~\cite{Tervo2014stochastic}. Значит, высокой фоновой \nt{НОРА} соответствует \emph{малое} действующее $g$, а пачке --- большое.

\section{Сколько искать}

Какое усиление лучше? Мы проверили это на модели. Сеть выбирает одно из двух действий по~\eqref{eq:softmax}; каждое приносит награду с какой-то вероятностью, и сеть оценивает её по наградам выбранного действия, как в главе~\ref{sec:dofa}. Время от времени мир меняется: обе вероятности разыгрываются заново. Меняться может и выбранное действие, и то, которое сеть давно не пробовала; о втором сеть может узнать только поиском. На рис.~\ref{fig:invertedU} показано, какую долю наилучшей возможной награды получает сеть при разных постоянных $g$.

\begin{figure}[t]
\centering
\begin{tikzpicture}[>=Stealth,x=1.2cm,y=1cm]
\draw[->,gray!70!black] (0,0) -- (7.6,0) node[below left,black] {\small $g$};
\draw[->,gray!70!black] (0,0) -- (0,4.4) node[above,black,font=\small] {доля наилучшей награды};
\foreach \x/\l in {0/0.5,1/1,2/2,3/4,4/8,5/16,6/32,7/64}{\draw[gray!60!black] (\x,-0.06) -- (\x,0.06); \node[below,font=\scriptsize] at (\x,0) {\l};}
\foreach \v/\y in {0.8/0.8,0.9/2.4,1.0/4.0}{\draw[gray!60!black] (-0.06,\y) -- (0.06,\y); \node[left,font=\scriptsize] at (0,\y) {\v};}
\draw[very thick,blue!60!black] plot[mark=*,mark size=1.3pt] coordinates {(0,0.368) (1,0.880) (2,1.728) (3,2.784) (4,3.456) (5,3.616) (6,3.488) (7,3.264)};
\draw[very thick,red!70!black] plot[mark=*,mark size=1.3pt] coordinates {(0,0.368) (1,0.672) (2,1.184) (3,1.840) (4,2.176) (5,1.920) (6,1.456) (7,1.104)};
\node[blue!60!black,font=\small,anchor=south] at (5,3.7) {спокойный мир};
\node[red!70!black,font=\small,anchor=north] at (5.1,1.25) {быстро меняющийся мир};
\draw[->,thick,blue!60!black] (5,3.25) -- (5,3.55);
\draw[->,thick,red!70!black] (4,1.8) -- (4,2.1);
\node[font=\small\itshape,gray!35!black,anchor=west] at (0.15,2.3) {ищет вслепую};
\node[font=\small\itshape,gray!35!black] at (6.45,2.55) {не замечает перемен};
\end{tikzpicture}
\caption{Доля наилучшей возможной награды при постоянном коэффициенте усиления $g$ (шкала по $g$ логарифмическая). Синяя кривая --- мир меняется в среднем раз в тысячу попыток, красная --- раз в двадцать. Стрелки отмечают лучшее $g$: в быстро меняющемся мире оно вдвое меньше. Среднее по 60 прогонам по 4000 попыток.}
\label{fig:invertedU}
\end{figure}

При $g=0.5$ сеть почти не пользуется знанием и получает около трёх четвертей возможного, при очень большом $g$ упускает перемены. Лучшее значение: $g=16$, когда мир меняется раз в тысячу попыток (доля $0.976$), и $g=8$, когда раз в двадцать ($0.886$).

\begin{keyblock}
\begin{proposition}[Перевёрнутое U]
\label{prop:explore}
Награда, которую собирает сеть, зависит от коэффициента усиления как перевёрнутое U: слишком малое $g$ тратит попытки на поиск, слишком большое не замечает перемен. Лучшее $g$ тем меньше, чем быстрее меняется мир.
\end{proposition}
\end{keyblock}

Перевёрнутое U наблюдается и в опытах: животные решают задачу лучше всего при умеренной фоновой активности \nt{НОРА}-нейронов с чёткими отрывистыми ответами, а при высокой фоновой активности отвлекаются и переключаются между задачами~\cite{AstonJones2005}. Это согласуется с различием режимов из предыдущего раздела: отрывистая пачка в момент решения даёт большое действующее $g$ как раз тогда, когда надо выбрать, а высокая фоновая активность без пачек усиливает всё подряд, в том числе посторонние входы, так что действующее $g$ для текущей задачи падает, --- это поиск.

\section{Кто крутит ручку}

Раз лучшее усиление зависит от того, как быстро меняется мир, его хорошо бы подстраивать. Но откуда \nt{НОРА}-нейронам знать, что мир изменился? Естественный признак --- \emph{неожиданность}: ошибки предсказания стали больше обычного. Важно различать два вида неопределённости. Если награда просто случайна, ошибки велики всегда, и это ожидаемо. Если же ошибки вдруг выросли по сравнению с их обычной величиной, значит, что-то поменялось. По одной из гипотез, именно эту неожиданную неопределённость и сообщает \nt{НОРА}, а ожидаемую --- \nt{АЦЕТ}~\cite{YuDayan2005}.

Что с этим сигналом делать? Можно понизить усиление и больше искать. Можно повысить скорость обучения, чтобы быстрее забыть устаревшие оценки: опыты показывают, что после резкой перемены люди учатся быстрее, и зрачок при этом расширяется~\cite{Nassar2012}; напомним, что зрачок --- показатель не только \nt{НОРА}, но и \nt{АЦЕТ}~\cite{Reimer2016}. А можно сделать и то и другое сразу: по ещё одной гипотезе отрывистая пачка \nt{НОРА} работает как \emph{прерывание} --- сигнал, по которому сеть сбрасывает текущее состояние и перестраивается под новую обстановку~\cite{BouretSara2005}, как общая линия прерывания процессора.

Честно скажем, чего мы не нашли. В модели мы пробовали оба простых способа подстройки: понижать $g$ или повышать скорость обучения, когда сглаженная ошибка превышает свою долговременную среднюю. В мире, который то подолгу стоит, то часто меняется, лучшие настройки обоих способов дали $0.926$--$0.927$ от наилучшей награды --- почти столько же, сколько лучшее постоянное $g$ ($0.925$ при $g=8$) или постоянная, но достаточно большая скорость обучения ($0.928$), а неудачные настройки --- меньше. Кривые на рис.~\ref{fig:invertedU} у вершины пологие, и в таком мире подстраивать почти нечего. Подстройка должна окупаться там, где разные режимы требуют очень разных настроек. Какой именно закон управления реализует \nt{НОРА}, пока не установлено.

\section{Итог}

\begin{table}[t]
\centering
\footnotesize
\setlength{\tabcolsep}{4pt}
\renewcommand{\arraystretch}{1.15}
\begin{tabular}{>{\raggedright}p{3.0cm}>{\raggedright}p{4.4cm}>{\raggedright\arraybackslash}p{6.0cm}}
\toprule
Что делает \nt{НОРА} & Как & Что известно \\
\midrule
меняет крутизну ответа нейронов & коэффициент усиления $g$ в~\eqref{eq:gain} & повышение отношения сигнала к шуму измерено; мультипликативная модель --- упрощение \\
переключает выбор & переводит сеть через $g\beta=1$ & следует из модели~\eqref{eq:wtagain}; прямых измерений порога нет \\
задаёт, сколько искать & $g$ --- обратная температура~\eqref{eq:softmax}; фон --- малое действующее $g$ (поиск), пачка --- большое (решение) & перевёрнутое U и два режима работы измерены; связь с поиском --- хорошо обоснованная гипотеза \\
сообщает о переменах & неожиданная неопределённость & зрачок и скорость обучения растут после перемен --- измерено; что это сигнал \nt{НОРА} --- гипотеза \\
прерывает и сбрасывает & отрывистая пачка на неожиданное событие & пачки на важные события измерены; «прерывание» --- гипотеза \\
\bottomrule
\end{tabular}
\caption{Что добавляет \nt{НОРА} к сети из \nt{ГЛУТ}, \nt{ГАМК} и \nt{ДОФА}.}
\label{tab:nora}
\end{table}

\nt{ДОФА} говорит сети, \emph{куда} менять веса. \nt{НОРА} не меняет ни одного веса, но решает, \emph{как} сеть пользуется тем, что уже знает (таблица~\ref{tab:nora}): одна ручка --- коэффициент усиления --- превращает нейрон из усилителя в компаратор, сеть выбора --- из «размышляющей» в «решившую», а выбор --- из поиска в использование. Как \nt{НОРА} узнаёт, насколько быстро меняется мир, --- открытый вопрос.

\section{Задачи}

\begin{exercise}[\lvlA]
\label{ex:08:1}
\emph{Одна ручка на весь мозг.}
(а)~По таблице~\ref{tab:types} оцените, сколько нейронов мозга приходится на один \nt{НОРА}-нейрон и сколько всего мест выброса у \nt{НОРА}-системы.
(б)~Пусть $\Delta u=1$, $\sigma_{\text{до}}=0.5$, $\sigma_{\text{после}}=4$. Вычислите $d'$ по~\eqref{eq:gainsnr} при $g=1$, $g=8$ и $g\to\infty$.
(в)~При каком $g$ достигается $90\,\%$ предельного $d'$? Покажите, что ответ зависит только от отношения $\sigma_{\text{после}}/\sigma_{\text{до}}$, и объясните смысл утверждения~\ref{prop:gainnoise} через это отношение.
\end{exercise}

\begin{exercise}[\lvlB]
\label{ex:08:2}
\emph{Выбор с коэффициентом усиления: вывод.} Система: $\tau\,\dd r_1/\dd t=-r_1+g[I_1-\beta r_2]_+$, $\tau\,\dd r_2/\dd t=-r_2+g[I_2-\beta r_1]_+$.
(а)~Выведите~\eqref{eq:wtagain}. Найдите собственные числа линеаризации в области, где работают обе группы, и время, за которое затухает разность $r_1-r_2$. Вычислите множитель усиления разницы и это время при $g\beta=0.5$ и $0.9$, $\tau=20\ms$.
(б)~Докажите, что состояние, где работают обе группы, существует только при $g\beta<(1-\varepsilon)/(1+\varepsilon)$ (при $\varepsilon>0$). Чему равна эта граница при $\varepsilon=0.05$?
(в)~Докажите утверждение из подписи к рис.~\ref{fig:pitchfork}: победа слабого входа устойчива при $g\beta>(1+\varepsilon)/(1-\varepsilon)$. Что происходит между двумя границами?
(г)~Объясните, почему у границы $g\beta=1$ сеть «размышляет» долго, и чем это полезно и вредно.
\end{exercise}

\begin{exercise}[\lvlA]
\label{ex:08:3}
\emph{Softmax из шума.}
(а)~Покажите, что при логистическом шуме $\eta$ с масштабом $\sigma$, $P(\eta<z)=1/(1+e^{-z/\sigma})$, вероятность выбора равна~\eqref{eq:softmax}.
(б)~Пусть каждая из $K$ групп получает свой независимый шум с распределением Гумбеля, $P(\eta_k<z)=\exp(-e^{-z/\sigma})$, и побеждает группа с наибольшим $gu_k+\eta_k$. Докажите, что $P(k)=e^{gu_k/\sigma}/\sum_le^{gu_l/\sigma}$, и выведите из этого случай $K=2$.
(в)~Вычислите $P(A)$ при $u_A-u_B=0.2$, $\sigma=1$ и $g=1$, $8$, $32$, и найдите $g$, при котором $P(A)=0.9$ при $u_A-u_B=0.1$.
(г)~Для гауссова шума $P(A)=\Phi\bigl(g(u_A-u_B)/s\bigr)$. Подберите $s$ так, чтобы наклоны в нуле совпали с логистической кривой, и найдите наибольшее расхождение кривых. Проверьте слова «кривая почти та же».
\end{exercise}

\begin{exercise}[\lvlB]
\label{ex:08:4}
\emph{Оценка лучшего усиления.} В модели рис.~\ref{fig:invertedU} вероятности наград двух действий после каждой перемены равномерно распределены на $[0,1]$.
(а)~Пусть сеть знает вероятности точно. Какую долю наилучшей награды она получает при $g=0.5$? Сравните с рис.~\ref{fig:invertedU}.
(б)~Пусть лучший вариант лучше худшего на $\Delta$. Сеть пробует худший вариант примерно раз в $e^{g\Delta}$ попыток. Потребуйте, чтобы это время было порядка среднего времени между переменами $1/h$, и получите $g^*\approx\ln(1/h)/\bar\Delta$. Найдите $\bar\Delta=\langle|p_A-p_B|\rangle$.
(в)~Вычислите оценку при $h=0.001$, $0.01$ и $0.05$ и сравните с лучшими $g$ модели. Чего оценка не учитывает?
\end{exercise}

\begin{exercise}[\lvlB]
\label{ex:08:5}
\emph{Моделирование: как лучшее усиление зависит от скорости перемен.} Работайте с копией \texttt{models/\allowbreak nora\_explore.py} в отдельной папке: скрипт пишет файлы \texttt{scan\_h*.txt} в текущую папку.
(а)~Для $h\in\{0.001,0.003,0.01,0.03,0.05,0.1,0.2\}$ и $g=2^{k/2}$, $k=-2,\dots,14$, найдите долю наилучшей награды (60 прогонов по 4000 попыток) и лучшее $g^*(h)$.
(б)~Постройте $g^*$ в зависимости от $\ln(1/h)$ вместе с оценкой задачи~\ref{ex:08:4}. В каком диапазоне оценка верна?
(в)~Объясните, почему при быстрых переменах $g^*$ перестаёт убывать. Указание: сравните время, за которое устаревает добытая поиском информация, со временем $1/\eta$, за которое правило $Q\leftarrow Q+\eta(R-Q)$ её усваивает, и оцените стационарный разброс $Q$ при $\eta=0.2$.
(г)~Как меняется высота вершины перевёрнутого U с $h$ и почему кривые у вершины пологие?
\end{exercise}

\begin{exercise}[\lvlB]
\label{ex:08:6}
\emph{Проектирование: прерывание для сети выбора.} Сеть задачи~\ref{ex:08:2} с $\beta=2$, $\tau=20\ms$, $g=1$ уже решила: $r_1=0.9$, $r_2=0$, а входы теперь $I_1=0.9$, $I_2=1.0$. По утверждению~\ref{prop:wta} старое решение держится, хотя второй вход сильнее. \nt{НОРА} посылает прерывание: на время $T_p$ усиление падает до $g_{\text{lo}}$, затем возвращается к $1$.
(а)~Покажите, что при $g=1$ исход определяется знаком $r_1-r_2-g(I_2-I_1)/(g\beta-1)$ в момент возврата усиления (пока работают обе группы).
(б)~При $g_{\text{lo}}=0$ найдите наименьшее $T_p$ аналитически.
(в)~Найдите наименьшее $T_p$ моделью (Эйлер, шаг $10$~мкс) при $g_{\text{lo}}=0$, $0.1$, $0.25$, $0.4$. Сформулируйте правило проектирования: сколько должна длиться пачка и насколько глубоко опускать усиление.
(г)~Почему прерывание через усиление лучше, чем прямое торможение победителя дополнительным \nt{ГАМК}-входом?
\end{exercise}

\begin{exercise}[\lvlC]
\label{ex:08:7}
\emph{Когда подстройка усиления окупается.} В разделе «Кто крутит ручку» подстройка $g$ по неожиданности почти ничего не дала. Выясните, где предел.
(а)~Оракул знает, спокойная сейчас эпоха или переменчивая, и ставит в каждой своё постоянное $g$. Найдите моделью лучшую пару и долю наилучшей награды в мире главы (эпохи по $1000$ попыток, $h=0.001$ и $0.05$). Сравните с лучшим постоянным $g$ и с подстройкой по неожиданности из \texttt{models/\allowbreak nora\_explore.py}. Какую часть возможного выигрыша подстройка уже забирает?
(б)~Спроектируйте мир, где оракул выигрывает у лучшего постоянного $g$ заметно больше. Используйте задачу~\ref{ex:08:4}: лучшее $g$ зависит и от $h$, и от $\bar\Delta$.
(в)~Проверьте в этом мире подстройку по неожиданности. Какую часть выигрыша оракула она забирает и почему?
(г)~Открытая часть. Какой сигнал должен получать \nt{НОРА}, чтобы приблизиться к оракулу? Сравните сигнал «сглаженная $|\delta|$ выше своей средней» с оценкой частоты перемен $h$ по последовательности наград. Какие измерения зрачка или активности \nt{НОРА}-нейронов могли бы различить эти две стратегии?
\end{exercise}
